\chapter{除法}\label{ch:g4:division}

三年级用乘法口诀和小的余数做除法
（见\cref{ch:g3:division}）；本章介绍完整的
笔算方法 --- 长除法 --- 它适用于任意的数，
按数级逐段进行计算。

\section{长除法}

\begin{method}[用一位数做除数的长除法]\label{met:g4:division:method}
计算 $749$ 除以 $6$：
\begin{enumerate}
  \item 取最左边的数字 $7$（百位）：$6$ 的 $1$ 倍不超过 $7$，
        商中写 $1$，减去
        $6$，余数为 $1$；
  \item \emph{落下}下一位数字 $4$：现在计算 $14$
        （十位）：$6 \times 2 = 12$；商写 $2$，余数 $2$；
  \item 落下 $9$：计算 $29$：$6 \times 4 = 24$；商写
        $4$，余数为 $5$；
  \item 没有剩余的数字了：商为 $124$，余数为 $5$。
\end{enumerate}
验算：$6 \times 124 + 5 = 744 + 5 = 749$，且 $5 < 6$。
\end{method}

\begin{omfigure}
\begin{tikzpicture}[scale=0.95]
  % digit columns at x = 0.2, 0.8, 1.4; every digit under its place
  \node[font=\normalsize] at (0.2,3.0) {$7$};
  \node[font=\normalsize] at (0.8,3.0) {$4$};
  \node[font=\normalsize] at (1.4,3.0) {$9$};
  % divisor bracket and quotient
  \draw[thick] (1.8,2.0) -- (1.8,3.35);
  \draw[thick] (1.8,2.62) -- (3.7,2.62);
  \node[font=\normalsize] at (2.35,3.0) {$6$};
  \node[font=\normalsize, omDef] at (2.35,2.2) {$1$};
  \node[font=\normalsize, omDef] at (2.95,2.2) {$2$};
  \node[font=\normalsize, omDef] at (3.55,2.2) {$4$};
  % step 1: 7 - 6 = 1
  \node[font=\normalsize, omThm] at (-0.35,2.4) {$-$};
  \node[font=\normalsize, omThm] at (0.2,2.4) {$6$};
  \draw (-0.55,2.1) -- (0.55,2.1);
  \node[font=\normalsize] at (0.2,1.8) {$1$};
  \node[font=\normalsize] at (0.8,1.8) {$4$};
  \draw[-Stealth, omExo, thin] (0.8,2.75) -- (0.8,2.05);
  % step 2: 14 - 12 = 2
  \node[font=\normalsize, omThm] at (-0.35,1.2) {$-$};
  \node[font=\normalsize, omThm] at (0.2,1.2) {$1$};
  \node[font=\normalsize, omThm] at (0.8,1.2) {$2$};
  \draw (-0.55,0.9) -- (1.15,0.9);
  \node[font=\normalsize] at (0.8,0.6) {$2$};
  \node[font=\normalsize] at (1.4,0.6) {$9$};
  \draw[-Stealth, omExo, thin] (1.4,2.75) -- (1.4,0.85);
  % step 3: 29 - 24 = 5
  \node[font=\normalsize, omThm] at (0.25,0.0) {$-$};
  \node[font=\normalsize, omThm] at (0.8,0.0) {$2$};
  \node[font=\normalsize, omThm] at (1.4,0.0) {$4$};
  \draw (0.05,-0.3) -- (1.75,-0.3);
  \node[font=\normalsize] at (1.4,-0.6) {$5$};
  % commentary
  \node[right, font=\small, omExo, align=left] at (4.5,1.4)
    {$7 - 6 = 1$，落下 $4$；\\
     $14 - 12 = 2$，落下 $9$；\\
     $29 - 24 = 5$：结束。\\[3pt]
     商为 $124$，余数为 $5$\\
     验算：$6 \times 124 + 5 = 749$};
\end{tikzpicture}

{\small 完整写出 $749$ 除以 $6$ 的长除法：每一步的减法都
列出，每个数字都留在自己的数位上，灰色
箭头表示被落下的数字。}
\end{omfigure}

\begin{example}[商中的零]\label{ex:g4:division:zero}
计算 $618$ 除以 $3$：百位，$6 \div 3 = 2$；十位，落下
$1$：$3$ 的 $0$ 倍不超过 $1$ --- \emph{一定要写} $0$！ ---
余数 $1$；个位，落下 $8$：$18 \div 3 = 6$。商为
$206$，余数 $0$。漏写中间的零（写成 $26$）
是最常见的错误；验算
$3 \times 26 = 78 \neq 618$ 能立即发现它。
\end{example}

\begin{example}[估算商的大小]\label{ex:g4:division:size}
在计算 $749$ 除以 $6$ 之前，先估算答案的范围：
$6 \times 100 = 600$，$6 \times 200 = 1\,200$，所以商在
$100$ 与 $200$ 之间 --- 商是三位数。这个一行的
估算能避免大多数数位放错的错误。
\end{example}

\section{看清题目在问什么}

\begin{method}[商、余数，还是商加一]\label{met:g4:division:interpret}
与\cref{met:g3:division:interpret}一样，由题意决定：
\begin{enumerate}
  \item “能装满的有多少\dots”或“每份有多少”：商；
  \item “还剩多少”：余数；
  \item “全部需要多少个容器”：商，
        若余数不为零，则再加一。
\end{enumerate}
\end{method}

\begin{example}\label{ex:g4:division:problem}
要把 $530$ 本书装箱，每箱装 $8$ 本。
\[
530 = 8 \times 66 + 2 .
\]
能装\emph{满}的箱子有几个？$66$。不在满箱中的书有几本？
$2$。要把\emph{全部}的书装下，需要几个箱子？$67$。
三个问题，一次除法。
\end{example}

\begin{example}[公平分钱]\label{ex:g4:division:money}
四个朋友平均分担一件礼物的 $92$ 元费用：
$92 \div 4 = 23$ 恰好除尽（余数为 $0$）。每人付 $23$。当
余数为零时，除法“正好除尽”，分配
十分公平。
\end{example}

\section{练习}

\begin{exercise}[$\star$]\label{exo:g4:division:1}
仿照\cref{ex:g4:division:size}，判断每个商
介于哪两个“整”数之间？再用竖式计算：
$96 \div 4$；\quad $87 \div 5$。
\end{exercise}

\begin{exercise}[$\star$]\label{exo:g4:division:2}
用竖式计算下列除法，并检验每一题：
$672 \div 4$；\quad $925 \div 7$；\quad $804 \div 6$。
\end{exercise}

\begin{exercise}[$\star$]\label{exo:g4:division:3}
注意零的处理（见\cref{ex:g4:division:zero}）：
$816 \div 4$；\quad $420 \div 7$；\quad $2\,512 \div 5$。
\end{exercise}

\begin{exercise}[$\star$]\label{exo:g4:division:4}
写出下列除法的检验等式（$a = b \times q + r$）：$85 \div 9$；
\quad $1\,000 \div 3$。
\end{exercise}

\begin{exercise}[$\star$]\label{exo:g4:division:5}
把 $156$ 个鸡蛋装入每箱 $6$ 个的箱子中。能装满的箱子有几个？
\end{exercise}

\begin{exercise}[$\star$]\label{exo:g4:division:6}
一根长 $250$ cm 的丝带，剪成每段 $8$ cm 的小段，能剪几段，
还剩多长？
\end{exercise}

\begin{exercise}[$\star$]\label{exo:g4:division:7}
$375$ 名学生乘坐每辆有 $52$ 个座位的大巴去看演出，需要几
辆大巴？（判断这是
\cref{met:g4:division:interpret}中的哪一种情况？用 $52$ 的倍数做除法：
$52 \times 7 = 364$。）
\end{exercise}

\begin{exercise}[$\star$]\label{exo:g4:division:8}
三个朋友平均分 $234$ 颗弹珠。每人分到几颗？
用乘法检验。
\end{exercise}

\begin{exercise}[$\star$]\label{exo:g4:division:9}
求被除数：用 $7$ 去除，商为 $58$，
余数为 $3$。
\end{exercise}

\begin{exercise}[$\star\star$]\label{exo:g4:division:10}
图书管理员要摆放 $438$ 本书，每层书架放 $9$ 本。
每个书柜配 $6$ 层书架。需要几层书架？需要买几
个书柜？（两次除法，都要
考虑“再多一个”的问题。）
\end{exercise}

\begin{exercise}[$\star\star$]\label{exo:g4:division:11}
用 $8$ 去除某数，商恰好等于余数。
所有可能的被除数是多少？请全部列出。（余数必须
小于 $8$。）
\end{exercise}
